% chapters/ch19_limitations.tex
% Chapter 19: Limitations and Future Work + High-Speed Adaptive Layer
% Source: NRMO_v5_updated.pdf, Chapter 19 (lines 8734-9146 of v5_layout.txt)

\chapter{Limitations and Future Work ---
High-Speed Adaptive Layer}
\label{ch:limitations}

This chapter combines the v2.0 limitations and future-work catalogue
with the formal definition of the high-speed adaptive layer
(NRMO REIN High-Speed Theory). It includes the engine architecture
and pseudocode that are the cognitive counterpart of the engineering
implementation in Part~X.

% =====================================================================
\section{Purpose and Layer Architecture}
\label{sec:limit-purpose}

\subsection{Subject of the High-Speed Theory}

NRMO Core specifies admissibility (the irreversible boundary). The
High-Speed Theory specifies how to act \emph{within} that boundary
when timescale collapses (e.g., market shock, incident, adversarial
encounter). The shorter timescale is the subject; the longer
timescale (Long-Term Layer / NRMO REIN) provides constraints.

\subsection{Formal Definitions}

\paragraph{State.}
$s \in S$, observation $x \in X$, action $a \in A$. State
transition: $s(t+1) \sim P(s(t), a(t))$.

\paragraph{Irreversible (Ruin).}

\begin{equation}
R \subset S, \qquad
s \in R \;\Rightarrow\; \forall u \ge t,\;\; s(u) \in R.
\end{equation}

\noindent
Purpose: Ruin avoidance.

% =====================================================================
\section{Layer Structure}
\label{sec:limit-layer-structure}

\subsection{NRMO Core (Veto)}
\label{sec:limit-nrmo-core}

\begin{equation}
\mathrm{adm}_N(x, a) \in \{\mathrm{YES},\, \mathrm{NO},\,
\mathrm{HOLD},\, \mathrm{SAFE\_EXIT}\}.
\end{equation}

\begin{itemize}
\item \textbf{YES}: action permissible.
\item \textbf{NO}: action prohibited.
\item \textbf{HOLD}: shutdown (defer for human review).
\item \textbf{SAFE\_EXIT}: emergency termination.
\end{itemize}

\subsection{Action Space and Execution}

\begin{equation}
A_H(x) \subset A; \qquad
\mathrm{execution} = \{a \in A_H(x) : \mathrm{adm}_N(x, a) =
\mathrm{YES}\}.
\end{equation}

% =====================================================================
\section{Five Engine Patterns}
\label{sec:limit-engine-patterns}

\subsection{A: Assumption Kill-Switch (Premise Monitor)}

Premise set $A_{\mathrm{ass}} = \{a_1, a_2, \ldots, a_n\}$;
kill-conditions $K_i: X \to \{0, 1\}$.

\begin{equation}
\exists i\; K_i(x) = 1
\;\Rightarrow\;
\mathrm{HOLD}\; \text{or}\; \mathrm{SAFE\_EXIT}.
\end{equation}

\noindent
Effect: any premise violation immediately demotes action space.

\subsection{B: Abort-First Design}

Abort condition set $B = \{b_1, \ldots, b_m\}$.

\begin{equation}
\exists b_j\; b_j(x, z) = 1 \;\Rightarrow\; \mathrm{SAFE\_EXIT}.
\end{equation}

\noindent
Effect: prioritises abort over continued action.

\subsection{C: Pre-Commit Trigger}

Trigger set $T = \{(C_k, A_k)\}$.

\begin{equation}
C_k(x) = 1 \;\Rightarrow\; a := A_k.
\end{equation}

\noindent
Effect: pre-committed action executes deterministically when
triggered.

\subsection{D: Explicit Expiry}

Each commitment $d = (\mathrm{content}, t_0, \Delta t)$ expires:

\begin{equation}
t > t_0 + \Delta t \;\Rightarrow\; d := \mathrm{expired}.
\end{equation}

\noindent
Effect: prevents indefinite commitment carry-forward.

\subsection{E: Fast / Slow Split Memory}

\begin{equation}
M_f(t) = \{m_i : t - t_i \le \tau\}; \qquad
M_s = \mathrm{archive}(M_f).
\end{equation}

\noindent
Effect: fast memory governs active operation; slow memory archives
for long-horizon reference. Their separation prevents short-horizon
transients from corrupting long-horizon constraints.

% =====================================================================
\section{Regime Profiles}
\label{sec:limit-regimes}

The five engine patterns are tuned per regime:

\paragraph{Market-Speed Regime.}
Premise: liquidity depth. Kill: gap/jump. Abort: drawdown.
Expiry: short. Notes: applicable to active trading.

\paragraph{Incident Regime.}
Premise: containment integrity. Kill: containment breach.
Abort: escalation. Expiry: 30--60 minutes. Notes: incident response.

\paragraph{Adversarial Regime.}
Premise: trust assumptions. Kill: trust breach. Abort: power
concession. Expiry: variable. Notes: adversarial encounter.

\paragraph{Long-Horizon Regime (Asset Allocation).}
Premise: macro alignment. Kill: regime change. Abort: irreversible
loss. Expiry: 6--12 months. Notes: long-horizon allocation.

% =====================================================================
\section{High-Speed Adaptive Layer: System Flow}
\label{sec:limit-system-flow}

\begin{lstlisting}[basicstyle=\ttfamily\scriptsize,frame=single]
+---------------------------------------+
|         HUMAN (Final Judgment)        |
+---------------------------------------+
                  |
+---------------------------------------+
| Long-Term Layer (life/daily scale)    |
| OUTER_SHELL (always-on)               |
+---------------------------------------+
                  | (on rupture)
+---------------------------------------+
| HUMIDITY_BUFFER (conditional)         |
+---------------------------------------+
                  | (on irreversible risk)
+---------------------------------------+
| Short-Term Layer: HIGH-SPEED          |
| ADAPTIVE                              |
| +-- Kill-Switch                       |
| +-- Abort-First                       |
| +-- Pre-Commit                        |
| +-- Expiry                            |
| +-- Fast Memory                       |
+---------------------------------------+
                  |
+---------------------------------------+
| APCSO (3-option generation)           |
+---------------------------------------+
                  |
+---------------------------------------+
| UNIFIED NRMO CORE (veto judgment)     |
+---------------------------------------+
                  |
+---------------------------------------+
| DECISION_SURVIVAL_ENGINE              |
+---------------------------------------+
                  |
+---------------------------------------+
| SLOW MEMORY (persistent)              |
+---------------------------------------+
\end{lstlisting}

% =====================================================================
\section{High-Speed Mode Pseudocode}
\label{sec:limit-pseudocode}

\begin{lstlisting}[language=Python,basicstyle=\ttfamily\small,frame=single]
HIGH_SPEED_MODE:
    # Kill-Switch evaluation
    for assumption in ASSUMPTIONS:
        if assumption.kill_condition(state):
            goto SAFE_EXIT

    # Pre-Commit evaluation
    for trigger in PRECOMMIT_TRIGGERS:
        if trigger.condition(state):
            action = trigger.action
            goto NRMO_CHECK

    # Abort-First evaluation
    if abort_condition():
        goto SAFE_EXIT

    # APCSO
    choice_set = apcso.generate_3_choices()

    # UNIFIED NRMO CORE
NRMO_CHECK:
    approved = []
    for action in choice_set:
        if unified_nrmo_core.check(action) == YES:
            approved.append(action)

    present_to_human(approved)

    # Expiry
    if action_executed:
        set_expiry(action, dt)
    fast_memory.append(action)

    # Slow memory archive
    if episode_finished:
        slow_memory.store(fast_memory.snapshot())
\end{lstlisting}

% =====================================================================
\section{OUTER\_SHELL}
\label{sec:limit-outer-shell}

Structure for maintaining daily continuity \emph{without} judgement.
Always present; no startup / shutdown concept.

\begin{lstlisting}[basicstyle=\ttfamily\small,frame=single]
PURPOSE: maintain_daily_continuity_without_decision
MODE:    always_on

ACCEPTS:
  - external_pressure
  - internal_urgency
  - ambiguity
  - no_event_state

OPERATIONS:
  - log_only
  - normalize
  - decay_over_time

FORBIDDEN:
  - interpretation
  - prioritization
  - decision_request
  - deadline_internalization

RUPTURE_CRITERIA:
  - daily_continuity_break:
      sleep / food / shelter / income x2 cycles
  - irreversible_risk:
      legal_or_physical_immediacy
  - external_coercion_with_third_party_harm

ON_RUPTURE:
  - escalate_to:  DECISION_SURVIVAL_ENGINE
  - bypass:       HUMIDITY_BUFFER (optional)
\end{lstlisting}

% =====================================================================
\section{HUMIDITY\_BUFFER}
\label{sec:limit-humidity-buffer}

Isolation layer to prevent emotional state from being converted to
judgement or tool.

\begin{lstlisting}[basicstyle=\ttfamily\small,frame=single]
PURPOSE:    preserve_affective_state
ACTIVATION: conditional

TRIGGERS:
  - urge_to_define
  - relational_density_high
  - silence_instability
  - observation_point_shift

ACTIONS:
  - hold_state
  - delay_response
  - schedule_remeasurement

FORBIDDEN:
  - evaluate_daily_continuity
  - consider_deadlines
  - assess_risk
\end{lstlisting}

% =====================================================================
\section{DECISION\_SURVIVAL\_ENGINE}
\label{sec:limit-decision-survival}

Final judgement engine for irreversibility avoidance and survival.

\begin{lstlisting}[basicstyle=\ttfamily\small,frame=single]
PURPOSE: decision_engine

INVARIANTS:
  - survival_first
  - irreversible_risk_avoidance
  - remeasurement_assumed
\end{lstlisting}

% =====================================================================
\section*{Connections to Other Parts}

\begin{itemize}
\item The five engine patterns A/B/C/D/E correspond to
implementation modules in Part~X: Kill-Switch / Abort-First / Pre-Commit /
Expiry / Memory split. The Python codebase
(\texttt{nrmo\_full/engine/}) maintains an analogous module structure
for the high-frequency civilisation simulator.
\item Regime profiles (market-speed / incident / adversarial /
long-horizon) correspond to the world families of Part~VIII (Normal,
LateStagnation, FastExpansionRace, Vulnerable, PlanetaryStress).
\item OUTER\_SHELL is the cognitive analogue of the always-on
governance layer in $\Omega$ Full (Part~IX).
\end{itemize}
